% Drawing 147/06-2 -- folder 147, batch 06, archive page 120 (PDF page 121).
% Transcription: courbe plane fermée qui se recoupe elle-même en une boucle
% Kind: closed plane curve with a self-intersection. Method 2 (schematic redraw from the facsimile). Status: ok.
% Notes: A closed plane curve crossing itself in a loop. Above it, a first small sketch of the same (a small loop with its two branches crossing) is drawn too.
% Coordinates: 1 unit = 1 mm = 1/100 of the crop width (bbox in page fractions: [0.514, 0.362, 0.655, 0.483]).
% Redrawn by Claude (claude-opus-5-5), issue #73, 2026-11-07; not checked by a human; not compiled (no TeX in the session).
% Rights: reproduces part of the work -- behind RENDER_TRANSCRIPTS, never published by itself.
\documentclass[tikz,border=3pt]{standalone}
\usepackage{amsmath,amssymb}
\usetikzlibrary{patterns,arrows.meta}
\tikzset{pencil/.style={black!45}, struck/.style={black!35}, mark/.style={circle,draw,inner sep=1.3pt}, pt/.style={circle,fill,inner sep=0.9pt}, opt/.style={circle,draw,fill=white,inner sep=1.1pt}}
\begin{document}
\begin{tikzpicture}[x=1mm,y=1mm,line cap=round,line join=round,font=\scriptsize]
\draw plot[smooth, tension=0.55] coordinates {(46,-9.8) (29.8,-21.3) (14.9,-27.7) (10.6,-24.5) (17,-20.2) (38.3,-20.2) (68.1,-26.8)};
\draw plot[smooth, tension=0.55] coordinates {(91.1,-33) (68.1,-38.3) (51.1,-49) (38.3,-68.1) (48.3,-85.1) (68.1,-83) (85.1,-70.2) (90,-57.5) (83,-42.6) (72.3,-39.4) (66,-44.7) (69.1,-72.4) (61.7,-95.8) (50.2,-117)};
\end{tikzpicture}
\end{document}
